\chapter*{LỜI MỞ ĐẦU}
\addcontentsline{toc}{chapter}{LỜI MỞ ĐẦU}

\section*{Lý do chọn đề tài}

Trong một chiến dịch marketing, người phụ trách thường phải làm đồng thời nhiều việc: xác định mục tiêu, lựa chọn kênh, chuẩn bị nội dung, lập lịch, theo dõi ngân sách và đọc các chỉ số sau khi đăng. Khi thông tin nằm rải rác ở nhiều bảng tính hoặc công cụ, việc kiểm tra trạng thái và tổng hợp kết quả mất nhiều thời gian. Nội dung do AI tạo ra có thể giúp giảm bớt công việc viết nháp, nhưng nếu không có giới hạn và người kiểm duyệt thì lại tạo thêm rủi ro về thông tin sai, quyền truy cập và trách nhiệm quyết định.

Đề tài “Hệ thống quản lý chiến dịch marketing có tích hợp trí tuệ nhân tạo” được lựa chọn để giải quyết hai nhu cầu trên trong một hệ thống nhỏ, phù hợp phạm vi học phần. Hệ thống tập trung vào quản lý dữ liệu và quy trình; AI đóng vai trò trợ lý gợi ý, tạo bản nháp và tóm tắt số liệu, không thay thế người duyệt.

\section*{Mục đích của tài liệu}

Đây là báo cáo dùng xuyên suốt toàn bộ môn AIA331. Ba bài kiểm tra thường xuyên và bài thi cuối được dùng như các tiêu chí chất lượng của cùng một sản phẩm, không tách thành các báo cáo riêng. Tài liệu trình bày:

\begin{enumerate}
  \item bối cảnh, người sử dụng, phạm vi và mục tiêu của hệ thống;
  \item yêu cầu chức năng, yêu cầu chất lượng và các quy tắc nghiệp vụ;
  \item kiến trúc, mô hình dữ liệu, quy trình phê duyệt và giao diện dự kiến;
  \item cách đặt AI vào đúng vị trí, giới hạn dữ liệu gửi đi và xử lý khi AI trả kết quả không hợp lệ;
  \item kế hoạch xây dựng, kiểm thử, đánh giá và triển khai.
\end{enumerate}

\section*{Phạm vi của bản tài liệu hiện tại}

Nhóm đang ở giai đoạn chuẩn bị tài liệu, chưa có mã nguồn, cơ sở dữ liệu chạy thật, dữ liệu marketing thật hay kết quả đánh giá mô hình. Vì vậy, các sơ đồ và bảng trong báo cáo là cơ sở để bắt đầu triển khai chứ không phải ảnh chụp của một sản phẩm đã hoàn thành. Đây là điểm cần lưu ý khi đọc các chương sau.

\section*{Bố cục báo cáo}

Chương 1 nêu bối cảnh, bài toán, mục tiêu và phạm vi. Chương 2 trình bày cơ sở tham khảo và nguyên tắc thiết kế. Chương 3 phân tích tác nhân, ca sử dụng và yêu cầu. Chương 4 mô tả kiến trúc; Chương 5 mô hình dữ liệu và quy trình. Chương 6 thiết kế các chức năng AI; Chương 7 đưa ra kế hoạch xây dựng, kiểm thử và triển khai. Chương 8 tổng kết và nêu hướng hoàn thiện. Phụ lục gồm bảng đối chiếu, prompt mẫu và DDL.
